% ---------------------------------------------------------------------------
\section{Problem, system, conventions and observables}\label{sec:setup}

\subsection{Question and objective}
The user asked (verbatim, translated where Chinese): \emph{``Theoretically, what is the highest
pressure or jetting velocity one could achieve from laser-induced cavitation? Please resolve this
completely.''} The handoff (\texttt{EXPERT\_HANDOFF.md}) further requires: (i) an independent audit of the
existing conditional bounds; (ii) reconciliation of the finite real-fluid PFC study; (iii) resolution of
activation, compressible/directional focusing, gel rupture and target impact; (iv) a defensible
\emph{conditional} maximum with validity, uncertainty and attainability; and (v) an experimental program
for whatever cannot be closed without data. The objective of this report is therefore not a single
number, but a hierarchy: rigorous ceilings that no process satisfying the stated hypotheses can
exceed, reduced-model predictions with verified and benchmarked components, and an explicit gap between them.

\subsection{Physical system}
One single-pulse event in which a nanosecond laser pulse is absorbed in or at a liquid
perfluorocarbon (PFC) droplet (perfluoro-\emph{n}-pentane, PFP, or perfluoro-\emph{n}-hexane, PFH)
embedded in an aqueous carrier (water, or a hydrogel), at ambient absolute pressure $\pinf$ and
temperature $T_0$. The PFC is superheated, nucleates, vaporises and expands as a cavity; the cavity
reaches a maximum volume and collapses. Near a boundary the collapse is asymmetric and can form a
liquid microjet. A receiver (target) surface may be loaded by the expansion, by the collapse pressure
wave, or by the jet. Laser-induced \emph{optical breakdown} (plasma formation) of the carrier is a
different activation route; it is discussed only to bound the general question (Section~\ref{sec:envelope})
and is excluded from the PFC benchmark by definition (hypothesis H0 below).

\paragraph{Coordinates and signs.} Spherical radius $r$ (outward positive) for spherical models;
for directed jets an axisymmetric cylinder system $(r,z)$ with the rigid receiver plane at $z=0$ and the
bubble centre initially at $z=h>0$. Wall (radial) velocity $\dot R>0$ means expansion; the jet speed
$U_{\mathrm j}$ is counted positive \emph{toward} the receiver ($-z$ direction). Pressures are
absolute unless called \emph{gauge} (relative to $\pinf$) or \emph{driving} (a stated difference). Time
$t=0$ at the end of the deposition stage unless stated.

\paragraph{Initial and boundary conditions (benchmark).} Liquid PFC at rest, stress-free relative to
$\pinf$ apart from Laplace pressure, temperature $T_0$; carrier at rest at $T_0$, $\pinf$; far-field
carrier quiescent at $\pinf$, $T_0$; rigid receiver with zero normal velocity (image condition);
free-field when no receiver is declared. The uniform, approximately isobaric activated state is an
adopted preparation (H1--H2), not a resolved optical deposition trajectory. The quantity
$\tau_{\mathrm h}$ used in CNT is the interval assigned to the hot-state nucleation hazard.
Its default 5\,ns is comparable to the 6\,ns acoustic transit time of a 6\,$\mu$m PFP droplet;
therefore it cannot by itself establish H2. A real short pulse may be stress confined
(Section~\ref{sec:isochoric}).

\subsection{Hypotheses used by the rigorous results}
\begin{description}[leftmargin=1.2cm,style=sameline]
\item[H0] Thermal activation: the absorbed optical energy appears as heat in the absorber/PFC; no
optical breakdown or plasma in the carrier.
\item[H1] Full activation: the absorbed heat raises the whole PFC droplet uniformly to the activation
state before significant motion (non-uniform, absorber-localised heating is discussed in Section~\ref{sec:validity}).
\item[H2] No stress confinement: during heating the liquid pressure stays near $\pinf$ plus Laplace
pressure (heating time long compared with the acoustic transit time of the droplet), unless the
isochoric route is explicitly invoked (Section~\ref{sec:isochoric}).
The reported activation EOS adopts the isobar $p_\ell=\pinf$, neglecting the finite-droplet
Laplace correction at preparation; retaining that correction would require recalculating the
spinodal, CNT and activated budget at the actual liquid pressure.
\item[H3] After deposition, the environment is a reservoir at $(T_0,\pinf)$; the system exchanges
heat only with it and receives no external useful work. Absorber/carrier/gel prestress and additional
heat sources are absent from the benchmark or must be included in its total activated exergy.
\item[H4] Continuum thermodynamics with the declared equations of state inside their validity
domains (Table~\ref{tab:ledger}).
\end{description}

\subsection{What is maximised: observable definitions (handoff priority 1)}\label{sec:observables}
The word ``maximum'' is meaningless until the observable is fixed. The following are distinct and
are never interchanged in this report.
\begin{description}[leftmargin=0.6cm]
\item[O1 Internal microjet speed $U_{\mathrm j}$.] Maximum speed of the liquid at the far pole of the
bubble, toward the receiver, before it touches the opposite bubble wall (definition of Supponen et
al.\ 2016, \S5.3, verified). Unit m/s.
\item[O2 Mass-resolved jet speed $U_{\mathrm j}(m_*)$.] The largest speed $U$ such that the jet mass
moving at least as fast as $U$ is at least $m_*$. Here $m_*>0$ is a specified participating
mass (kg); $u_{\mathrm{jet}}$ is the local speed toward the receiver (m/s), $\dd m$ a jet mass element (kg),
$m_{\mathrm j}(\ge U)$ the mass faster than $U$ (kg), and $\sup$ the least upper bound over nonnegative speeds:
\begin{equation}
U_{\mathrm j}(m_*) \coloneqq \sup\{\,U\ge 0:\ m_{\mathrm j}(\ge U)\ge m_*\,\},\qquad
m_{\mathrm j}(\ge U)=\int_{u_{\mathrm{jet}}\ge U}\dd m .\label{eq:Um}
\end{equation}
The application chooses $m_*$, e.g.\ the mass needed to load an ink or register in a sensor.
This gives a mass-qualified observable instead of an unconstrained point-tip maximum.
\item[O3 Arrival speed $U_{\mathrm{arr}}$.] Normal speed of liquid reaching the receiver plane.
\item[O4 Source pressure $p_b$.] Absolute pressure inside the cavity during expansion.
\item[O5 Collapse-wave pressure $p_{\mathrm{ff}}(d)$.] Peak pressure increment of the outgoing wave
at distance $d$ from the cavity centre, in the far field of the spherical model.
\item[O6 First-contact receiver pressure $p_{\mathrm c}$.] Pressure increment at the receiver at the
instant of normal liquid contact (shock/Hugoniot matching), lasting until lateral release.
\item[O7 Averaged receiver pressure $\bar p(A_{\mathrm o},\tau_{\mathrm o})$.] Area/time average over a
declared observer area $A_{\mathrm o}$ and duration $\tau_{\mathrm o}$ (Eq.~E32 of the limits report).
\end{description}
Initial optical-breakdown shocks (route excluded by H0) are not part of O4--O7 for the PFC benchmark.

\subsection{Benchmark architectures and admissible design set}\label{sec:benchmark}
No calibrated specimen was supplied (Section~\ref{sec:ledger}). Following the handoff
(``choose and clearly label a defensible benchmark''), two benchmarks are declared.
\begin{description}[leftmargin=0.6cm]
\item[B1 (liquid carrier, rigid receiver).] One PFC droplet in water at stand-off $h$ from a flat rigid
receiver; observables O1--O6.
\item[B2 (gel-embedded).] The same droplet in a neo-Hookean/Kelvin--Voigt hydrogel; observables O4--O5
and the rupture gate (Section~\ref{sec:gel}); directed jets in gel are not solved.
\end{description}
For the design set, $d_0=2a_0$ is initial diameter (m); $\tau_{\mathrm h}$ the hot-state hazard
interval (s); $\gamma=h/\Rm$ the dimensionless stand-off with initial centre height $h$ and
maximum free-field cavity radius $\Rm$ (m); $\alpha_{\mathrm c}\in[0,1]$ the dimensionless
phase-transfer coefficient; $x_{\mathrm{air}}$ the ambient gas volume per PFC liquid volume
(dimensionless); and $G$ the gel shear modulus (Pa). The broad admissible control set is
\begin{equation}
\begin{aligned}
\mathcal D=\Bigl\{&\text{species}\in\{\mathrm{PFP},\mathrm{PFH}\},\ d_0\in[3,12]\,\mu\mathrm m,\
\tau_{\mathrm h}\in[1\,\mathrm{ns},1\,\mu\mathrm s],\\
&\gamma=h/\Rm\in[1.25,4],\ \alpha_{\mathrm c}\in[0,1],\ x_{\mathrm{air}}\in[0,0.3],\ G\in[0,20]\,\mathrm{kPa}\Bigr\},
\end{aligned}\label{eq:design}
\end{equation}
Fixed: $T_0=293.15$\,K,
$\pinf=101.325$\,kPa, water carrier properties at $T_0$, full activation of the droplet (uniform
heating to its kinetic superheat limit). The lower end $\gamma\ge1.25$ excludes wall contact during
growth, where the potential-flow method used here is not valid (Section~\ref{sec:jet}).
The jet runs sample the 6\,$\mu$m, 5\,ns, water, air-free B1 slice at five stand-offs
and four content closures; they do not optimize every continuous coordinate of $\mathcal D$.

\subsection{Notation}
Symbols are defined where first used; the recurring ones are collected here (units SI).
\begin{center}\small
\begin{longtable}{>{\raggedright}p{0.17\textwidth}p{0.62\textwidth}p{0.13\textwidth}}
\toprule Symbol & Meaning & Unit\\ \midrule\endhead
$T_0,\ \pinf$ & environment (dead-state) temperature and absolute pressure & K, Pa\\
$\Eabs$ & optical energy absorbed in the system during the pulse & J\\
$\mathcal U,\mathcal S,V$ & internal energy, entropy, volume of the system & J, J/K, m$^3$\\
$u,s,v,h_{\mathrm F}$ & specific internal energy, entropy, volume, PFC enthalpy & J/kg, J/(kg\,K), m$^3$/kg, J/kg\\
$\Phi,\ \psi$ & exergy (availability) of the system relative to $(T_0,\pinf)$; specific exergy & J, J/kg\\
$W_{\mathrm u}$ & useful work: total work minus work against $\pinf$ & J\\
$\mF,\ a_0,\ d_0$ & PFC mass, initial droplet radius, diameter & kg, m, m\\
$\rho_{\mathrm F},\ \rho_{\mathrm w},\ c_{\mathrm w}$ & PFC liquid density; carrier density and sound speed & kg/m$^3$, kg/m$^3$, m/s\\
$\psat(T)$ & saturation pressure from the reference EOS & Pa\\
$\Tsp(p)$ & liquid-spinodal temperature at liquid pressure $p$ & K\\
$\Tk$ & kinetic (CNT) superheat limit & K\\
$\sigma_{\mathrm F},\ \sigma$ & PFC liquid/vapour tension; cavity-wall tension in the carrier & N/m\\
$R,\ \dot R,\ \Rm$ & cavity radius, wall velocity, maximum radius & m, m/s, m\\
$p_b,\ p_L$ & cavity pressure; carrier pressure at the wall & Pa\\
$\alpha_{\mathrm c}$ & condensation/evaporation coefficient (Hertz--Knudsen--Schrage) & --\\
$h,\ \gamma,\ \zeta$ & stand-off, $\gamma=h/\Rm$, anisotropy parameter & m, --, --\\
$\Delta p_{\mathrm d}$ & collapse driving pressure used for scaling & Pa\\
$U_{\mathrm j},\ U^*_{\mathrm j}$ & microjet speed; $U^*_{\mathrm j}=U_{\mathrm j}/\sqrt{\Delta p_{\mathrm d}/\rho_{\mathrm w}}$ & m/s, --\\
$\phi,\ q$ & velocity potential; its normal derivative (scaled in Section 6) & m$^2$/s, m/s\\
$\Gf,\ \Hf$ & Green's function with wall image and its normal-derivative kernel & 1/m, 1/m$^2$\\
$\KE(m),\ \EE(m)$ & complete elliptic integrals of the first/second kind, parameter $m$ & --\\
$\mathrm{We}$ & Weber number $\Delta p_{\mathrm d}\Rm/\sigma$ & --\\
$n_{\mathrm T},\ B_{\mathrm T}$ & Tait exponent and pressure for water & --, Pa\\
$G,\ G_{\mathrm c}$ & gel shear modulus; fracture energy & Pa, J/m$^2$\\
\bottomrule
\end{longtable}
\end{center}
Equation numbers carry the prefix X to avoid confusion with the E01--E32 numbering of the limits
report and of the finite-PFC study; those are always cited with their document
(``Limits E19'', ``Finite-PFC E27'').
